\section{Evaluation Model and Assumptions}\label{sec:methodology}
\subsection{Scheduler Panel}
\label{sec:panel}

The PRIMARY panel comprises 11 policies, executed alongside 2 additional
\texttt{STYLE\_APPROXIMATION} policies used only for a mandatory robustness
check (13 policies executed in total): a classical arrival-order baseline
(\texttt{fifo}, also the load-calibration reference policy), a
deadline-aware policy (\texttt{edf}), a slack/urgency-aware policy
(\texttt{least\_laxity\_first}), a service-length-aware policy
(\texttt{estimated\_service\_time\_first}), a fairness-aware policy
(\texttt{weighted\_fair\_share}), a KV/memory-pressure-aware policy
(\texttt{kv\_constrained\_online}), an admission-control/SLO-guard policy
(\texttt{admission\_control}), and four historically named simulator implementations of
external mechanisms: vLLM-style token-budget continuous batching in both
its pre-chunked-prefill and chunked-prefill forms
(\texttt{vllm\_faithful}, \texttt{vllm\_chunked\_prefill\_faithful}),
Sarathi-style chunked prefill (\texttt{sarathi\_faithful}), and an
SLO-aware scheduling policy in the SLAI/RAD family
(\texttt{slai\_faithful}). The two \texttt{STYLE\_APPROXIMATION} robustness
policies are a token-budget-proxy style approximation
(\texttt{vllm\_style\_token\_budget}) and an SLO-guard-proxy style
approximation (\texttt{scorpio\_style\_slo\_guard}); two further implementations requiring disaggregated/cross-instance simulator
semantics (\texttt{distserve\_faithful}, \texttt{llumnix\_faithful}) are
designated for a secondary stratum, never pooled with the single-GPU
colocated PRIMARY leaderboard, and one candidate
(\texttt{apt\_serve\_faithful}) is excluded entirely as
\texttt{scaffolding\_only} -- not a validated reimplementation at the time
of this freeze.

A policy enters the PRIMARY panel only if, at the time it was added to the
comparability audit -- frozen a full day before the discriminability
pilot (\S\ref{sec:stage0}) executed -- it satisfies an outcome-independent
selection rule: its mechanism predates this protocol's freeze, it
occupies a mechanism family not already saturated by an existing panel
member, it was judged to have adequate implementation fidelity and passed its own unit
tests, it is simulator-semantically compatible with the single-GPU
colocated execution model, and it is independently reproducible from a
cited reference. No policy was added to or removed from the panel in
response to the pilot's outcome, and mechanism families it identified as
under-implemented (preemptive/multilevel scheduling in the FastServe
sense, JITServe, LMetric) were deliberately left as documented future-work
candidates rather than fast-tracked into the panel after seeing the
pilot's BurstGPT collapse.

\subsection{Execution Model}
\label{sec:execution-model}

Each (window, policy, operating-region) cell is executed as a deterministic
discrete-event simulation. The entities are requests, one or more simulated
GPU resources, a waiting queue, and policy-controlled active batches. Each
request carries an arrival time, prompt-token count, target output-token
count, priority weight, and, where required by a policy or metric, an
SLO/deadline value. Simulator state includes the pending-arrival list,
waiting queue, active requests per GPU, request phase and progress
(prefill remaining, decoded tokens, first-token time), current KV-token
footprint, per-step active-batch and utilization histories, completed and
dropped requests, and observational mechanism telemetry
(\S\ref{sec:telemetry}).

Time advances in fixed serving steps. At each step, all requests with
arrival time less than or equal to the current simulation time are enqueued;
the simulator exposes an observable state to the scheduling policy; the
policy returns an action specifying admissions and, for policies that support
them, preemption, swapping, migration, held decodes, or prefill-control
decisions; the simulator validates the action against the resource model and
applies it; active requests then advance by one service step; completed
requests are removed and recorded; and per-step utilization, batch, KV, and
contention histories are updated. The primary benchmark uses the colocated
single-GPU-compatible stratum of this model; disaggregated or cross-instance
semantics exist for secondary policies but are not pooled with the PRIMARY
leaderboard (\S\ref{sec:panel}).

The resource model includes active-sequence capacity, batch/token-budget
constraints where a faithful policy uses them, and KV-memory accounting.
KV occupancy is tracked as demand relative to the configured nominal
\texttt{max\_kv\_tokens}; only KV-aware policies enforce this limit at
admission, so occupancy above 1.0 is interpreted as demand exceeding nominal
capacity rather than as an instrumentation failure. In the default
backward-compatible mode, prefill is instantaneous and each active decoding
request produces one output token per step. In the prefill-aware mode used
by faithful serving-style policies, requests may first consume prefill
budget before decoding begins; decoding requests produce one token per step;
chunked-prefill policies may limit the prefill work admitted into a step;
and prefill/decode contention and token-budget saturation are recorded as
diagnostics. Admission-control and SLO-aware policies may decline otherwise
waiting requests or hold decode work according to their policy semantics;
SLO/deadline outcomes are then evaluated through the metric contract
(\S\ref{sec:metrics}).

Simulation terminates when all arrivals have been processed and no request
remains waiting, active, migrating, or relocating, or when the configured
post-arrival drain limit is reached. Requests left in the waiting or
transfer states at termination are treated as dropped. Metrics are computed
from completed requests, dropped requests, the full arrival population, and
the per-step histories; no completion-conditioned metric is imputed when no
request completes (\S\ref{sec:metrics}). Given identical requests,
configuration, policy, and operating region, the simulator consumes no
randomness and produces bit-identical outputs; this was verified by a
repeated-run determinism test. This determinism is what permits the paired
design of \S\ref{sec:unit}: two policies' outcomes share the same input window without independent
simulation noise. Differences also include policy-dependent service modes,
resource enforcement and information access; they cannot be attributed to
scheduling order alone.

\begin{table}[t]
\centering
\caption{Main simulation-model scope and assumptions used in the PRIMARY
ranking-portability study.}
\label{tab:model-assumptions}
\small
\begin{tabular}{p{0.25\linewidth}p{0.68\linewidth}}
\toprule
Model component & Treatment in this study \\
\midrule
Arrivals & Trace-driven request arrivals from frozen 200-request windows. \\
Service progression & Fixed-step serving model; decode emits one token per active decoding request per step. \\
Prefill/decode & Instant prefill in legacy mode; explicit prefill budget and contention diagnostics in prefill-aware modes. \\
Resource limits & Active-sequence, token-budget, and KV-demand constraints as implemented by the configured policy/model. \\
Scheduling decisions & Policies observe simulator state and decide admission, ordering, and supported control actions. \\
Metrics & Computed from arrivals, completions, drops, deadlines, and per-step histories; undefined conditional metrics are excluded, not imputed. \\
Omitted hardware effects & No claim of absolute latency or throughput fidelity to physical GPUs; kernel, networking, allocator, and hardware-noise effects are abstracted in the numerical campaign. \\
\bottomrule
\end{tabular}
\end{table}


\subsection{Historical fidelity labels and execution limits}\label{sec:fidelity}\label{sec:vv-structure}
The frozen registry distinguishes repository-native classical policies,
simulator proxies, style approximations and a class named
\texttt{FAITHFUL\_EXTERNAL}. We retain identifiers such as
\texttt{slai\_faithful} solely as immutable data keys. These names and their
historical unit-test gates do not certify native execution equivalence.
Policy-specific shadow progress and memory state need not equal the shared
execution state. Per-request token intentions are not always carried as
exact native runner plans; the resource model is not a native KV allocator.

Source-verified follow-up work found different final-prefill sampling,
admission and recomputation semantics in pinned vLLM 0.27.1. Native
scheduling can advance optimistic progress before runner acknowledgment;
recomputation retains previously emitted output for replay. Historical
simulator eviction can instead reset output and first-token progress.
These differences are plausible sources of comparative bias, but no
measured ranking correction or reversal direction follows from discovering
them. We neither repair historical outputs nor relabel them as v2 results.

The numerical campaign remains useful as a controlled comparison of the
specified implementations. Determinism verifies repeat execution within
that model. It does not establish absolute hardware timing accuracy,
mechanism equivalence to production systems, or physical portability of
relative rankings.
\subsection{Mechanism Telemetry}
\label{sec:telemetry}

Every cell records seven groups of simulator telemetry: queue depth;
active-batch saturation; prefill/decode contention; KV occupancy;
admission activations; preemption/reorder events; and token-budget
saturation. Means and maxima supplement several groups. A zero can mean
that an operation did not occur or that the configured model does not
represent that mechanism. These are simulator observations, not GPU
counters. The descriptor association analysis in
Section~\ref{sec:telemetry-explanation} does not establish a causal
explanation of the observed reversals.

\subsection{Implementation identity and comparison scope}
\begin{table}[t]\centering
\caption{Frozen panel interpretation. Registry labels are historical
identifiers; none of these rows certifies native serving equivalence.}\label{tab:panel}
\begin{tabular}{p{0.40\linewidth}p{0.51\linewidth}}\toprule
Policy or group & Historical comparison role\\\midrule
\texttt{fifo} & Arrival-order baseline and FIFO calibration reference\\
\texttt{edf}, \texttt{least\_laxity\_first} & Deadline and slack/urgency ordering\\
\texttt{estimated\_service\_time\_first} & Service-length-aware ordering under the simulator's information model\\
\texttt{weighted\_fair\_share} & Priority-weighted fairness discipline\\
\texttt{kv\_constrained\_online} & Simulator KV-pressure proxy\\
\texttt{admission\_control} & Admission/SLO guard\\
\texttt{vllm\_faithful} & Historical continuous-batching implementation\\
\texttt{vllm\_chunked\_prefill\_faithful} & Historical chunked-prefill implementation\\
\texttt{sarathi\_faithful} & Historical Sarathi-inspired implementation\\
\texttt{slai\_faithful} & Historical SLAI-inspired implementation\\
\texttt{vllm\_style\_token\_budget}, \texttt{scorpio\_style\_slo\_guard} & Two style approximations outside the primary ranking panel\\\bottomrule
\end{tabular}
\end{table}
A policy label is not a complete treatment definition. The shared simulator
service model, admission constraints, policy-local shadow state and input
information together determine the observed outcome. When a policy uses
request-length information that is available in a trace-driven replay,
that information is part of the historical experiment; it must not be
represented as an online estimate validated on a production engine.
Comparability therefore means controlled replay of the frozen inputs and
implementations, not identical internal information or resource enforcement
across every named mechanism.

The execution schema records configuration, source, window, policy and
repetition identities alongside simulator and manifest hashes. Reproducing
a table requires those identities as well as its numerical values. A
future corrected simulator or common-native adaptation must receive a
separate version and evidence namespace. It cannot replace a policy's old
implementation while retaining the old results as if they described the
new semantics. That version boundary permits an honest historical
comparison without freezing known limitations into future serving work.
